\section{The same geometric divided Frobenius induces the original Cartier map}

The source is the original exterior cotangent sheaf on the actual
Frobenius twist, and the primary target is the existing pushed de Rham
complex of an actual smooth Witt lift of a declared pure relative
dimension. The dimension is an explicit hypothesis of this standard
Deligne--Illusie input. The arbitrary-smooth intrinsic complex gives a
subsequent compatible target through its genuine local chart exchange. The affine Cartier theorem is used for the existing
canonical divided differential of the same chosen lift; no new affine
Cartier calculation is introduced.

\begin{lemma}[A whole map out of a single complex retains its original cycle class]
\label{lem:di-single-map-original-cycle-class}
\uses{lem:native-module-complex-every-cycle-class}
\lean{EnvelopeProperty.DI.LeadingCartierCycleTransport.singleClassMap, EnvelopeProperty.DI.LeadingCartierCycleTransport.singleClassMap_eq_liftCycles, EnvelopeProperty.DI.LeadingCartierCycleTransport.singleClassMap_cycleClass}
\leanok
Let $E$ be a nonnegative cochain complex of modules and
$\phi:M\to E^q$ a whole morphism whose image is closed. The map on
homology induced by the actual cochain map $M[-q]\to E$, after the
canonical identification $H^q(M[-q])=M$, sends every $x\in M$ to the
class of the original cycle $\phi(x)$ in this same $E$.
\end{lemma}
\begin{proof}
The map $\phi$ factors uniquely through the actual cycles $Z^q(E)$.
Naturality of cycles for the cochain map out of $M[-q]$ identifies its
cycles map with this factor, preceded by the canonical cycles
identification of the single complex. Naturality of the canonical
projection from cycles to homology then identifies its induced class
map with this cycle factor followed by $Z^q(E)\to H^q(E)$.
For any original cycle with value $\phi(x)$, the cycle inclusion sends
both the factor's value and that cycle to $\phi(x)$. The cycle inclusion
is a monomorphism, so the two cycles agree. Applying the original
homology projection proves the claimed exact class formula.
\end{proof}

\begin{lemma}[Canonical evaluation retains the original single-source homology map]
\label{lem:di-single-source-canonical-evaluation}
\uses{lem:di-single-map-original-cycle-class,thm:cartier-geometric-affine-homology-sections}
\lean{EnvelopeProperty.DI.SingleFunctorHomology.sectionComponent, EnvelopeProperty.DI.SingleFunctorHomology.sectionComponent_closed, EnvelopeProperty.DI.SingleFunctorHomology.singleClassMap_evaluation}
\leanok
Let $M[-q]\to K$ be an actual cochain map of original module sheaves,
and let $W$ be an original open. Evaluate its degree-$q$ component,
preceded by the original degree-$q$ single-source unit. The induced
single-source class map of this actual section map, followed by the
canonical comparison from homology of sections to sections of
homology, equals evaluation of the original sheaf homology map,
preceded by the original single-source homology identification.
If $W$ is affine and the terms of $K$ are quasicoherent, this same
comparison is the existing canonical affine homology-sections
isomorphism. No source exactness or desired class equality is assumed.
\end{lemma}
\begin{proof}
Evaluation preserves kernels and zero morphisms. Let $L$ be the unique
cycle factor of the evaluated original degree map. Compose $L$ with
the preserved-kernel comparison to evaluation of the original target
cycles. This agrees with evaluation of the original single-source
cycle identification followed by the original cycles map: composing
both maps with evaluation of the cycle inclusion gives the identical
original degree component. The preserved cycle inclusion is a
monomorphism, so the two whole cycle maps agree.

The canonical comparison from homology of sections takes the class of
an original cycle to evaluation of its original homology class. Apply
this identity to the preceding whole cycle equality. Naturality of the
projection from cycles to homology and the original single-source
cycles/homology identities give the asserted whole homology equality.
On an affine open with quasicoherent target terms, the existing affine
homology-sections isomorphism has precisely this canonical forward map,
which proves the final assertion.
\end{proof}

\begin{theorem}[The whole original geometric degree-one class is Cartier]
\label{thm:di-geometric-original-leading-cartier-one}
\uses{lem:di-single-map-original-cycle-class,lem:di-single-source-canonical-evaluation,thm:cartier-geometric-affine-homology-sections,thm:cartier-geometric-local-native-cartier}
\lean{EnvelopeProperty.DI.GeometricLeadingCartier.xiSections, EnvelopeProperty.DI.GeometricLeadingCartier.xiSections_apply, EnvelopeProperty.DI.GeometricLeadingCartier.xiSections_closed, EnvelopeProperty.DI.GeometricLeadingCartier.xiCochain, EnvelopeProperty.DI.GeometricLeadingCartier.xiCochain_f_one, EnvelopeProperty.DI.GeometricLeadingCartier.xiHomologyClass, EnvelopeProperty.DI.GeometricLeadingCartier.xiClassSections, EnvelopeProperty.DI.GeometricLeadingCartier.xiClassSections_eq_canonicalComparison, EnvelopeProperty.DI.GeometricLeadingCartier.xiClassSections_eq_inverseCartier}
\leanok
Let $f:X\to\operatorname{Spec}k$ be smooth of declared pure relative
dimension $n$, and let $p$ be prime with $\operatorname{char}k=p$.
On an original affine open $U$, put $B=\Gamma(U,\mathcal O_X)$ and let
$T=B\otimes_{k,F_k}k$ be its genuine native twist, with the original
geometric relative-Frobenius map $\beta:T\to B$. Fix an original etale
tuple of length $n$ in $B$. Let $C\to k$ and $R\to T$ be actual
reductions modulo $p$, and $S\to B$ an actual flat reduction modulo $p$.
Require $S$ formally smooth over $C$, the original coefficient actions
to equal these reduction actions, and an actual $C$-algebra lift
$F:R\to S$ whose whole reduction square is $\beta$. These are original
lifting data, not a Cartier-class conclusion. On a genuine affine chart
of the actual smooth Witt lift they are the internally produced
reductions and the same chosen lift.

Let $\xi$ be the canonical divided differential of this same $F$.
Transport its source by the original exterior source unit and the
actual coefficient-preserving twist comparison, and its target by the
original exterior target unit and canonical relative-Frobenius
preimage comparison. The homology map of the actual resulting local
representative, followed by the canonical affine homology-to-sections
comparison for the same original global complex, is the original
geometric inverse Cartier map on the whole original source module.
No desired Cartier equality or quasi-isomorphism is supplied as an input.
\end{theorem}
\begin{proof}
Let $T$ be the genuine native twist of the original residue chart
coordinate ring. The same actual source and target reductions and the
same chosen lift instantiate the canonical native divided one-form
map. Its image is closed by the existing native closedness theorem.
For an original geometric source section, apply the actual source unit
inverse, then the coefficient-preserving native cotangent comparison.
Call the resulting native cotangent element $\omega$.

The literal geometric target-unit formula identifies the image under
$\xi$ with the image of the original native cycle under the whole
original chart complex isomorphism. Scalar restriction preserves this
same cycle and its original class, and the actual complex isomorphism
preserves its class by functoriality of homology. The preceding lemma
therefore identifies the induced local homology map with the class of
that same native divided one-form, transported by these exact original
comparisons.

The existing affine degree-one Cartier theorem identifies this native
class with inverse Cartier applied to the degree-one exterior element
corresponding to $\omega$. The source comparisons in this formula are
exactly the original source exterior unit and the original twist
cotangent comparison, so the source is the original geometric section.
Finally apply the canonical affine homology-to-sections isomorphism.
Its whole forward map is the original canonical evaluation comparison,
so the target is the homology sheaf of the same intrinsic global
complex. These are precisely the defining source, target and scalar
comparisons of the original geometric inverse Cartier map. Equality
holds for every source section, hence for the entire module morphism.
\end{proof}

\begin{corollary}[The same actual Witt chart map has the original geometric Cartier class]
\label{cor:di-actual-witt-chart-leading-cartier}
\uses{thm:di-geometric-original-leading-cartier-one,lem:di-single-source-canonical-evaluation,thm:cartier-geometric-global-original-chart-formula}
\lean{EnvelopeProperty.DI.ActualWittLeadingCartier.sectionMap, EnvelopeProperty.DI.ActualWittLeadingCartier.xiSectionMap_eq, EnvelopeProperty.DI.ActualWittLeadingCartier.xiSectionMap_closed, EnvelopeProperty.DI.ActualWittLeadingCartier.representative, EnvelopeProperty.DI.ActualWittLeadingCartier.representative_f_one, EnvelopeProperty.DI.ActualWittLeadingCartier.classSections, EnvelopeProperty.DI.ActualWittLeadingCartier.classSections_eq_homologyRepresentative, EnvelopeProperty.DI.ActualWittLeadingCartier.classSections_eq_inverseCartier}
\leanok
Let $\widetilde X\to\operatorname{Spec}W_2(k)$ be an actual smooth lift
of declared pure relative dimension $n$, with $k$ perfect of
characteristic $p$. Fix an original affine lift chart, an honest residue
etale tuple, and the same selected relative Frobenius lift $F$ on that
chart. The original geometric section morphism $\Xi_F$ obtained from
this actual lift, with its original source exterior unit and target
exterior unit, induces the restriction of the original global
geometric inverse Cartier map on the entire degree-one source.
For an actual sheaf cochain representative with this original degree
component, canonical evaluation identifies its sheaf homology map with
this same section class map.
\end{corollary}
\begin{proof}
Use the actual geometric residue-section equivalence of this original
affine lift chart. The Witt coefficient reduction, its source reduction
along this equivalence, and its flat target reduction are genuine
reductions of the required kind. The whole lift-square identity for
the same $F$ is the original residue Frobenius identity transported
along this equivalence. The coefficient actions agree with these
reductions by their definitions, and formal smoothness is inherited
from the actual smooth Witt chart.

The source degree-one exterior identification, the actual twist section
comparison, the native divided differential, and the original target
exterior unit in the preceding theorem are precisely the four
successive factors defining $\Xi_F$. Their whole morphisms therefore
agree. The preceding theorem computes the class of this same map as
the original local geometric inverse Cartier map. The original global
Cartier isomorphism has exactly this whole local component on the same
honest chart, proving the first assertion. For an actual sheaf cochain
representative, apply canonical single-source evaluation to its original
degree component. It gives evaluation of the same original sheaf
homology map, proving the second assertion without assuming any class
equality or quasi-isomorphism.
\end{proof}

\begin{lemma}[Zero-source homology retains the same single-degree representative]
\label{lem:di-zero-source-single-class}
\uses{lem:di-single-map-original-cycle-class}
\lean{EnvelopeProperty.DI.ZeroSourceHomology.homologySelfIso, EnvelopeProperty.DI.ZeroSourceHomology.component_closed, EnvelopeProperty.DI.ZeroSourceHomology.homologyMap_eq_singleClassMap}
\leanok
Let $G$ be a complex whose differentials vanish and let $\Phi:G\to E$
be an actual whole cochain map. In degree $q$, identify $H^q(G)$ with
$G^q$ by its original cycle inclusion and cycle-to-homology projection.
The resulting homology map of $\Phi$ is exactly the homology map of the
actual single-degree representative of the same component $\Phi^q$.
\end{lemma}
\begin{proof}
Since the source differential vanishes, the original cycle inclusion
$Z^q(G)\to G^q$ and the original projection $Z^q(G)\to H^q(G)$ are
isomorphisms. The chain-map identity makes the image of $\Phi^q$
closed. Naturality of the cycle inclusion identifies the original
source cycle identification followed by the cycles map of $\Phi$
with the unique factorization of $\Phi^q$ through $Z^q(E)$: compose
both maps with its monomorphic cycle inclusion to obtain $\Phi^q$.
Naturality of the cycle-to-homology projection then says that the
normalized original homology map of $\Phi$ is this cycle factor
followed by the original target homology projection. This is also the
class map of its same single-degree representative, by the original
single-source class identity.
\end{proof}

\begin{theorem}[The original local exterior representative has the original Cartier class]
\label{thm:di-geometric-original-leading-cartier-exterior}
\uses{thm:di-geometric-original-leading-cartier-one,lem:di-zero-source-single-class,thm:cartier-geometric-affine-homology-sections,thm:cartier-geometric-global-original-chart-formula}
\lean{EnvelopeProperty.DI.GeometricLeadingCartierExterior.exteriorSections, EnvelopeProperty.DI.GeometricLeadingCartierExterior.exteriorSections_apply, EnvelopeProperty.DI.GeometricLeadingCartierExterior.exteriorSections_sourceWedge, EnvelopeProperty.DI.GeometricLeadingCartierExterior.exteriorCochain, EnvelopeProperty.DI.GeometricLeadingCartierExterior.exteriorCochain_f, EnvelopeProperty.DI.GeometricLeadingCartierExterior.exteriorHomologyClass, EnvelopeProperty.DI.GeometricLeadingCartierExterior.exteriorClassSections, EnvelopeProperty.DI.GeometricLeadingCartierExterior.exteriorClassSections_eq_canonicalComparison, EnvelopeProperty.DI.GeometricLeadingCartierExterior.exteriorClassSections_eq_inverseCartier, EnvelopeProperty.DI.GeometricLeadingCartierExterior.exteriorClassSections_eq_globalInverseCartier}
\leanok
For every exterior degree $q$, the exterior map constructed from the
same original $\xi$, evaluated through the same original source and
target units, induces the original degree-$q$ geometric inverse Cartier
map on the whole original source exterior module.
\end{theorem}
\begin{proof}
On each original affine chart, the whole exterior source comparison
sends every original wedge to the wedge of its actual native cotangent
images. The whole target exterior comparison sends the exterior
product of the same divided one-forms to its original geometric
exterior product. The canonical native all-degree Cartier theorem
therefore applies to the same exterior cochain map, with no new choice
of lift or representative. It identifies the native induced homology
map in degree $q$ with native inverse Cartier on that whole exterior
module. Scalar-restriction homology, the whole original chart complex
isomorphism, and the canonical affine homology-to-sections comparison
preserve the same cycles and their classes. Applying these comparisons
gives the stated whole original geometric map, in every degree,
including degree zero and zero coordinate rings.
\end{proof}

\begin{corollary}[The same actual Witt exterior map has the original global Cartier class]
\label{cor:di-actual-witt-exterior-leading-cartier}
\uses{cor:di-actual-witt-chart-leading-cartier,thm:di-geometric-original-leading-cartier-exterior}
\lean{EnvelopeProperty.DI.ActualWittExteriorCartier.exteriorSections, EnvelopeProperty.DI.ActualWittExteriorCartier.exteriorSections_sourceWedge, EnvelopeProperty.DI.ActualWittExteriorCartier.representative, EnvelopeProperty.DI.ActualWittExteriorCartier.representative_f, EnvelopeProperty.DI.ActualWittExteriorCartier.classSections, EnvelopeProperty.DI.ActualWittExteriorCartier.classSections_eq_homologyRepresentative, EnvelopeProperty.DI.ActualWittExteriorCartier.classSections_eq_inverseCartier, EnvelopeProperty.DI.ActualWittExteriorCartier.exteriorSections_one_eq_xiSectionMap}
\leanok
On an actual smooth Witt lift of declared pure dimension, the same
actual residue-section equivalence, original reductions and selected
lift produce the whole original exterior section map in every degree.
It takes every original source wedge to the original target-unit image
of the wedge of the same actual divided one-forms. Its degree-one map
is the same original geometric $\Xi_F$. Its actual native homology,
followed by the canonical original affine homology-sections map, is the
whole component of the original global geometric Cartier isomorphism.
\end{corollary}
\begin{proof}
Use exactly the internally produced reductions and whole lift square
of the degree-one Witt chart comparison. The all-degree original-unit
transport theorem applies to those same data. Its source-wedge formula
is the native exterior formula for the same canonical divided
one-forms; the actual residue-section construction defines these
one-forms using precisely those same reductions and lift. In degree
one the native exterior map followed by the native degree-one exterior
identification is the original native divided differential. The
original source and target units therefore give the same $\Xi_F$ on
the whole original module. The native cochain representative has this
same original component by its single-source construction, and its
actual homology is the class map in the all-degree transport theorem.
That theorem identifies the whole map with the original global
Cartier component in every degree.
\end{proof}
